\documentclass[12pt]{article}
\usepackage[utf8]{inputenc}
\usepackage[T1]{fontenc}
\usepackage{hyperref}
\title{Malay Language Resources for K-12 Curriculum Development}
\author{Compiled by Claude Sonnet 5.5 (claude-sonnet-5-5) for Andr\'as Kornai}
\date{}
\begin{document}
\maketitle

\section{Introduction}
Malay (ISO 639-3 zlm; Glottocode mala1479) is an Austronesian language of wider communication in Malaysia, Indonesia and nearby countries. The prompt suggests it may no longer have native speakers, but it is widely spoken as a first language and is a school subject in Malaysia. The searches here were narrow, focused on Malaysian dictionaries, and did not examine school curricula.

\section{Primary Educational Resources}
No textbooks or school curricula were retrieved by the search. The national language authority in Malaysia, Dewan Bahasa dan Pustaka (DBP), is the main publisher of reference works.

\section{Linguistic and Reference Materials}
\begin{itemize}
\item Wikipedia, ``Kamus Dewan'': the standard DBP monolingual dictionary, compiled by Teuku Iskandar. \url{https://en.wikipedia.org/wiki/Kamus_Dewan}
\item SAYS, article on the online Kamus Dewan Perdana from DBP (free access during the launch period, subscription planned). \url{https://says.com/my/tech/kamus-dewan-perdana-online}
\item Sinar Harian, ``Budi dan Kamus Dewan Perdana''. \url{https://www.sinarharian.com.my/article/106374/budi-dan-kamus-dewan-perdana}
\item UKM technical report on a Kamus Dewan related system (PDF). \url{https://ftsm.ukm.my/v5/file/research/technicalreport/PTA-FTSM-2020-104.pdf}
\item Kawah Buku, Kamus Melayu-Khmer Dewan, a bilingual dictionary. \url{https://kawahbuku.com/titles/kamus-melayu-khmer-dewan-%e1%9e%9c%e1%9e%85%e1%9e%93%e1%9e%b6%e1%9e%93%e1%9e%bb%e1%9e%80%e1%9f%92%e1%9e%9a%e1%9e%98-%e1%9e%98%e1%9f%89%e1%9e%b6%e1%9e%a1%e1%9e%b6%e1%9e%99%e1%9e%bc-%e1%9e%81%e1%9f%92}
\end{itemize}

\section{Religious and Community Resources}
No religious or community resources were retrieved.

\section{Digital and Audio Resources}
The search also mentioned an older free online dictionary, Pusat Rujukan Persuratan Melayu (prpm.dbp.gov.my). Only the host name appeared in the article text, not a full URL in the search results, so it is not listed as a resource.

\section{Curriculum Development Context}
Malay is a major language with extensive existing curricula, so the task is mainly to collect and organise materials rather than to revive anything. This dossier is thin and a fuller survey of Malaysian and Indonesian school materials would be needed.
\end{document}